\chapter{1945:Artturi Virtanen}

\label{chap:chemistry1945}

The Nobel Prize in Chemistry for the year 1945 was awarded to Artturi Virtanen.

\textbf{Motivation:}

for his research and inventions in agricultural and nutrition chemistry, especially for his fodder preservation method

\section{Artturi Virtanen}

\subsection*{Early Life and Education}

Artturi Virtanen was born on 1895-01-15 in Helsinki, Finland.

\subsection*{Career and Major Research}

Virtanen studied chemistry at the University of Helsinki and spent his career in agricultural and food chemistry. He worked at the research laboratory of the Valio cooperative, the Finnish Butter Export Association, where he studied the preservation of green fodder. In the 1920s he developed the AIV method, named after his initials, in which fresh forage is acidified with dilute hydrochloric or sulphuric acid to a pH of about 3 to 4 to stop spoilage fermentation. He later became professor at the Helsinki University of Technology.

\subsection*{Nobel Prize-winning Work}

The work for which Artturi Virtanen was awarded the Nobel Prize in Chemistry in 1945 was described by the Nobel Committee as: "for his research and inventions in agricultural and nutrition chemistry, especially for his fodder preservation method".

The AIV method preserved silage with far lower losses of nutrients than traditional ensiling or haymaking, while keeping the fodder palatable to livestock. The process was introduced into farm practice in Finland in 1929.

\subsection*{Significance and Impact}

Virtanen's method solved a practical problem of northern agriculture, where the growing season is short and fodder had to be stored for winter feeding. His demonstration that a low pH stops the degradation of proteins and vitamins became the basis of modern silage science. The method spread to other countries with cold, humid climates.

\subsection*{Later Career and Honors}

Virtanen continued research in biochemistry, including work on vitamins and on the fixation of nitrogen by the root nodules of leguminous plants. He remained at the Helsinki University of Technology for most of his career. He died in 1973.

\subsection*{Legacy}

Virtanen's name survives in the AIV process, and the principle of acid-based fodder preservation remains in worldwide use in ensiling practice.

\subsection*{Selected Publications}

\begin{itemize}

\item \emph{The A.I.V. method of preserving fresh fodder}, Empire Journal of Experimental Agriculture, 1933.

\end{itemize}

\clearpage

\section*{Chapter Summary}

The 1945 prize recognized a method for preserving cattle fodder that cut the heavy losses of nutrients suffered in traditional silage and haymaking. Artturi Virtanen's AIV process applied dilute mineral acid to fresh forage to prevent fermentative spoilage, and it came into wide use in Finland and other northern countries.
